% ============================================================
%  Préambule commun — Arbitrage Hub (Horacle Capital)
%  Police : Times (mathptmx), format article de recherche
% ============================================================
\usepackage[utf8]{inputenc}
\usepackage[T1]{fontenc}
\usepackage{mathptmx}            % Times New Roman (texte + maths)
\usepackage{courier}            % police à chasse fixe vectorielle
\usepackage{amsmath,amssymb,amsthm,bm}
\usepackage[a4paper,margin=2.4cm,headheight=14pt]{geometry}
\usepackage{microtype}
\usepackage{graphicx}
\usepackage{booktabs,tabularx,multirow,array,longtable}
\usepackage{enumitem}
\usepackage[dvipsnames,table]{xcolor}
\usepackage{fancyhdr}
\usepackage{titlesec}
\usepackage[font=small,labelfont=bf]{caption}
\usepackage[numbers,sort&compress]{natbib}
\usepackage{listings}
\usepackage[most]{tcolorbox}
\usepackage{needspace}
\usepackage{abstract}
\usepackage[hidelinks,unicode]{hyperref}

\frenchspacing
\setlength{\parindent}{1.2em}
\setlength{\parskip}{0.35em}
\linespread{1.06}
\setlist{itemsep=0.15em,topsep=0.3em}
\hyphenpenalty=600
\emergencystretch=2em

% ---- Libellés en français (pas de babel-french sur la machine de build)
\renewcommand{\contentsname}{Table des matières}
\renewcommand{\refname}{Références}
\renewcommand{\abstractname}{Résumé}
\renewcommand{\figurename}{Figure}
\renewcommand{\tablename}{Tableau}
\renewcommand{\listfigurename}{Liste des figures}
\newcommand{\og}{\guillemotleft\,}
\newcommand{\fg}{\,\guillemotright}

% ---- Titres de sections
\titleformat{\section}{\normalfont\large\bfseries}{\thesection.}{0.6em}{}
\titleformat{\subsection}{\normalfont\normalsize\bfseries}{\thesubsection.}{0.5em}{}
\titleformat{\subsubsection}{\normalfont\normalsize\itshape}{\thesubsubsection.}{0.5em}{}

% ---- En-têtes
\pagestyle{fancy}
\fancyhf{}
\renewcommand{\headrulewidth}{0.4pt}
\fancyfoot[C]{\small\thepage}

% ---- Théorèmes
\theoremstyle{definition}
\newtheorem{definition}{Définition}[section]
\newtheorem{example}{Exemple}[section]
\newtheorem{exercise}{Exercice}[section]
\theoremstyle{plain}
\newtheorem{proposition}{Proposition}[section]
\newtheorem{theorem}{Théorème}[section]
\newtheorem{lemma}{Lemme}[section]
\theoremstyle{remark}
\newtheorem{remark}{Remarque}[section]
\renewcommand{\proofname}{Démonstration}

% ---- Encadrés pédagogiques
\tcbset{boxrule=0.5pt,arc=1pt,left=6pt,right=6pt,top=4pt,bottom=4pt,fonttitle=\bfseries\small}
\newtcolorbox{intuition}[1][]{colback=Blue!3,colframe=Blue!55!black,title={Image mentale#1}}
\newtcolorbox{whiteboard}[1][]{colback=Goldenrod!6,colframe=Goldenrod!70!black,title={Au tableau blanc#1}}
\newtcolorbox{pitfall}[1][]{colback=Red!3,colframe=Red!60!black,title={Piège / limite#1}}
\newtcolorbox{toread}[1][]{colback=Gray!5,colframe=Gray!70!black,title={Pour aller plus loin — quoi lire#1}}
\newtcolorbox{codebox}[1][]{colback=Green!3,colframe=Green!45!black,title={Vérification en Python#1}}

% ---- Code
\lstset{basicstyle=\ttfamily\footnotesize,breaklines=true,columns=fullflexible,
  keywordstyle=\color{Blue!70!black},commentstyle=\color{Gray},stringstyle=\color{Maroon},
  frame=single,rulecolor=\color{Gray!40},xleftmargin=0.5em,xrightmargin=0.5em,
  literate={é}{{\'e}}1 {è}{{\`e}}1 {à}{{\`a}}1 {ç}{{\c c}}1 {ê}{{\^e}}1 {ô}{{\^o}}1 {î}{{\^i}}1 {ù}{{\`u}}1 {É}{{\'E}}1 {—}{{---}}1 {→}{{$\to$}}1}

% ---- Notations
\newcommand{\E}{\mathbb{E}}
\newcommand{\Var}{\operatorname{Var}}
\newcommand{\Cov}{\operatorname{Cov}}
\newcommand{\Corr}{\operatorname{Corr}}
\newcommand{\R}{\mathbb{R}}
\newcommand{\Prob}{\mathbb{P}}
\newcommand{\iid}{\overset{\text{i.i.d.}}{\sim}}
\newcommand{\N}{\mathcal{N}}
\newcommand{\dd}{\mathrm{d}}
